\section{Gestión integral, recursos y viabilidad}
\subsection{Alcance del proyecto y del servicio}
El producto comprende el CMS, clientes de consulta, tres terminales, conectividad, alimentación y contenidos iniciales. El proyecto comprende diagnóstico, especificación, adquisición, integración, pruebas, transferencia y evaluación del piloto. La operación posterior mantendrá contenidos, equipos y servicio dentro de las condiciones aceptadas.

\cuadro{tab:entregables}{Entregables del ciclo completo}{@{}P{0.7cm}P{3.4cm}Y@{}}{\textbf{ID}&\textbf{Entregable}&\textbf{Condición de aceptación}}{
E1 & Diagnóstico y decisión & Necesidad, interesados, línea base y alternativa documentados\\*
E2 & Diseño y especificación & Requisitos, ubicaciones, arquitectura, materiales y configuración conciliados\\
E3 & Adquisición y preparación & Equipos y licencias conformes; plataforma y contenidos preparados\\
E4 & Servicio integrado & Casos funcionales, de seguridad, continuidad y energía superados\\
E5 & Transferencia & Manuales, inventario, formación, respaldo y responsables de operación\\
E6 & Evaluación y cierre & Indicadores del piloto, costos observados y decisión de continuidad o ampliación\\
}
El control de cada entregable relacionará requisito, responsable, evidencia y aceptación. La línea base de alcance se aprobará al cerrar el diseño y se modificará mediante el procedimiento de cambios.

\subsection{Equipo formulador y coordinación institucional}
\cuadro{tab:equipo}{Responsabilidades de formulación e integración}{@{}P{3.7cm}Y P{3.2cm}@{}}{\textbf{Integrante}&\textbf{Responsabilidad}&\textbf{Revisión cruzada}}{
Cabeza Huillca, Flavio Antony, 211265 & Problema, usuarios, contenidos, coordinación y presentación & Presupuesto y aceptación\\*
Muñoz Castillo, George Ivanok, 204800 & Requisitos, clientes, plataforma, seguridad e integración & Configuración energética y riesgos\\
Fernández Baca Castro, Lucian Neptalí, 184197 & Costos, planificación, balance de formulación y control documental & Fuentes y coherencia del conjunto\\
}
La ejecución requiere además responsables institucionales de contenidos y TI, administración para las adquisiciones y especialistas para los trabajos eléctricos, fotovoltaicos y estructurales. La dirección de escuela coordinará estos interlocutores y su disponibilidad.

El esfuerzo previsto para servicios de implementación es de 408 horas-persona: 40 de ingeniería, 96 de instalación eléctrica, FV, red y videovigilancia, 100 de CMS, 60 de interfaz y contenido, 48 de pruebas, 24 de documentación y formación y 40 de coordinación. Los costos y las horas de operación se presupuestan por separado. El cronograma considera esperas de suministro y actividades concurrentes, además de esfuerzo de trabajo.

\subsection{Viabilidad y condiciones de inversión}
\cuadro{tab:viabilidad}{Evaluación integrada y evidencia de cierre}{@{}P{2.5cm}Y Y@{}}{\textbf{Dimensión}&\textbf{Fundamento}&\textbf{Evidencia de cierre}}{
Servicio y usuarios & Encuesta descriptiva y proceso editorial propuesto & Línea base local y ensayo de consulta\\*
Técnica TI & CMS y clientes disponibles; arquitectura modular & Prueba del conjunto y soporte de plataforma\\
Energética & Balance reproducible y recurso regional de referencia & Carga real, sombras, cubierta y configuración revisadas\\
Económica & Escenarios comparables, inversión y operación & Ofertas por conjunto y financiamiento de inversión y operación\\
Operativa & Responsables, capacitación y mantenimiento previstos & Designación, disponibilidad y transferencia aceptadas\\
Institucional & Uso del espacio, publicación y adquisición coordinados & Autorizaciones, licencias y procedimiento definidos\\
}
Se propone avanzar mediante puntos de decisión: confirmar necesidad, aceptar diseño, autorizar adquisición, aceptar el servicio integrado y evaluar su expansión. Cada punto requiere evidencia y una decisión registrada.

\subsection{Marco institucional, datos y ciclo ambiental}
La adquisición se tramitará mediante la instancia competente de la universidad, considerando la Ley N.° 32069, su reglamento y las modificaciones recogidas en el compendio oficial del MEF \cite{contrataciones}. Las especificaciones describirán desempeño y equivalencia; las marcas se utilizan para comparar opciones. La entidad determinará el procedimiento, la certificación presupuestal y los documentos de contratación aplicables.

La consulta pública utilizará avisos autorizados. Los contactos de colaboración se publicarán con autorización y finalidad definidas; las cuentas administrativas se limitarán a las funciones necesarias. El tratamiento de datos personales se revisará conforme a la Ley N.° 29733 y su reglamento \cite{datos}. La integración con sistemas institucionales requerirá autorización del titular y acuerdo de acceso.

Se mantendrán inventario, garantías, licencias, registros de mantenimiento y procedimientos de baja. Equipos electrónicos, baterías y módulos al final de su vida útil se entregarán a los sistemas de manejo que correspondan, con trazabilidad patrimonial y ambiental \cite{raee}. Se favorecerán componentes reparables, repuestos accesibles y sustitución del cliente sin descartar una pantalla funcional.

\subsection{Gestión de riesgos}
La probabilidad (P) y el impacto (I) se valoran de 1 a 3. Para P: 1, exposición limitada con control comprobado; 2, evento posible con control parcial; 3, exposición relevante sin control probado. Para I: 1, corrección local; 2, retrabajo o reprogramación; 3, afectación de seguridad, aceptación, viabilidad o un hito crítico. P$\times$I de 6 a 9 corresponde a prioridad alta, de 3 a 4 media y de 1 a 2 baja.

La valoración inicial refleja hipótesis de exposición y se revisará al obtener datos. Cuando falte fundamento se marcará por valorar. La respuesta se considera eficaz cuando su evidencia reduce la exposición o limita el efecto a una condición aceptada.

\cuadro{tab:riesgos1}{Riesgos de servicio, plataforma y energía}{@{}P{0.8cm}P{2.7cm}P{0.9cm}Y@{}}{\textbf{ID}&\textbf{Evento}&\textbf{P/I}&\textbf{Señal, respuesta y revisión}}{
R1 & Aviso vencido o sin autorización & 3/3 & Señal: aviso sin fuente o fecha. Flavio: exigir aprobación y caducidad; revisar muestras semanales de publicaciones\\*
R2 & Cliente pierde contenido tras reinicio & 3/3 & Señal: arranque sin contenido local. George: ensayo sin red y recuperación; aceptar sólo configuración persistente\\
R3 & Plataforma sin soporte o licencia adecuada & 2/3 & Señal: versión sin actualización o licencia pendiente. George: matriz de soporte y derechos; verificar antes de adquirir y trimestralmente en operación\\
R4 & Energía o autonomía insuficientes & 3/3 & Señal: carga o producción fuera del escenario. Lucian y especialista: medir, ajustar horario y redimensionar; revisar diseño y ensayo de corte\\
R5 & Incompatibilidad ambiental o eléctrica & 2/3 & Señal: límite de altitud, temperatura, corriente o BMS incumplido. Especialista: excluir combinación y demostrar variante compatible antes de instalar\\
}

\cuadro{tab:riesgos2}{Riesgos económicos, de instalación y operación}{@{}P{0.8cm}P{2.7cm}P{0.9cm}Y@{}}{\textbf{ID}&\textbf{Evento}&\textbf{P/I}&\textbf{Señal, respuesta y revisión}}{
R6 & Oferta incompleta o suministro tardío & 3/3 & Señal: exclusión esencial, vencimiento o cambio de stock. Lucian: comparar conjuntos, actualizar reserva y secuencia; revisión en cada oferta\\*
R7 & Anclaje, cubierta o circuito inseguros & 2/3 & Señal: diseño o prueba sin conformidad. Infraestructura: revisión competente y suspensión del montaje hasta corregir\\
R8 & Acceso indebido o datos expuestos & 2/3 & Señal: escape del kiosco o publicación no autorizada. George y TI: permisos, segmentación y pruebas negativas; revisión antes del servicio y tras cambios\\
R9 & Contenido sin actualización o soporte & 2/3 & Señal: avisos fuera de vigencia o incidencia sin responsable. Flavio: relevo, capacitación y cola de atención; revisión editorial semanal\\
R10 & Uso menor al esperado o expansión prematura & 2/3 & Señal: consultas fallidas y costo elevado por servicio útil. Equipo y escuela: piloto medido, ajuste de interfaz y evaluación por edificio\\
R11 & Daño, hurto o manipulación de un terminal & 2/3 & Señal: golpes, puertos abiertos o retiro de piezas. Infraestructura y TI: anclaje, acceso con llave, cámara por terminal y revisión diaria; exportar la grabación sólo con solicitud autorizada\\
}
El registro incorporará fecha, fundamento, estado, acción, evidencia y riesgo residual. Se revisará semanalmente durante la ejecución y en cada cambio de equipo, ubicación, presupuesto o requisito.

\subsection{Seguimiento y control de cambios}
El tablero Kanban organiza Por hacer, En curso y Listo, con un responsable y revisión por actividad. El tablero de formulación INFORMA contiene dos épicas y 16 tareas asignadas. La ejecución del piloto se desagregará a partir de la EDT y del cronograma de este informe, tras la autorización de recursos.

Cada solicitud de cambio registrará motivo, requisito afectado, impacto en costo, plazo, calidad, operación y riesgo. La coordinación propondrá aceptar, diferir o rechazar; el responsable institucional resolverá los cambios de inversión o alcance. La configuración aprobada, el presupuesto y los casos de aceptación se actualizarán en la misma versión.
